% 6. Descripción de la Competencia; 6.1. Estructura del Torneo; 12.1. Para Othello.
\subsection{Competencia externa reportada}
\textit{Guía: 6. Descripción de la Competencia; 6.1. Estructura del Torneo; 12.1. Para Othello.}
El equipo Ispis reportó cinco encuentros contra dos grupos, utilizando la evaluación \texttt{fija}, altura 4 y Minimax con poda Alpha-Beta. Se desconocen las configuraciones rivales, fecha, límite de tiempo y aperturas; no se registraron tiempos ni nodos. La tabla~\ref{tab:othello_externa} conserva los resultados reportados. Las tres capturas contra Juan Pablo respaldan los contadores por color; la asignación de equipos procede del usuario. Se conservan sin edición en la carpeta de evidencia externa y en un fragmento opcional de apéndice.
% 6.1. Estructura del Torneo; 12.1. Para Othello; 16. Formato de Tablas y Figuras.
\begin{table}[htbp]
\centering\small
\caption{Competencia externa reportada. I: Ispis; JP: equipo Juan Pablo; D: equipo Daniela; N/B: negras/blancas. Inició según colores reportados (negras). P5: terminalidad no confirmada.}
\label{tab:othello_externa}
\begin{tabular}{rlcccrrl}
\toprule
Partida & Rival & Color I & Color rival & Inició & Fichas I & Rival & Ganador \\
\midrule
1 & JP & N & B & I & 25 & 39 & JP \\
2 & JP & B & N & JP & 31 & 33 & JP \\
3 & JP & N & B & I & 25 & 39 & JP \\
4 & D & N & B & I & 49 & 15 & I \\
5 & D & B & N & D & 4 & 59 & D \\
\bottomrule
\end{tabular}
\end{table}


Ispis obtuvo una victoria y cuatro derrotas (20\% descriptivo), sin empates, un punto con puntuación 1/0.5/0 y diferencia media de $-10.20$ fichas. Contra Juan Pablo perdió los tres encuentros, con diferencia media $-10$; contra Daniela obtuvo una victoria y una derrota, con diferencia media $-10.50$. No se conoce un encuentro Juan Pablo--Daniela y no se presenta una clasificación completa de todos contra todos.

\textit{Guía: 3.1. Reglas Fundamentales; 13. Métodos Estadísticos Requeridos; 23. Política de Integridad Académica.}
Las partidas 1 y 3 comparten marcador y configuración propia; las imágenes muestran distribuciones visibles distintas, pero no se conocen sus inicios ni trayectorias. Se conservan ambas sin asumir independencia. La partida 5 suma 63 fichas: podría haber terminado legalmente con una casilla vacía, pero no se confirmó que ambos jugadores carecieran de movimientos. Se identifica como resultado reportado de terminalidad no verificada. Sin esa partida, los otros cuatro resultados darían una victoria, tres derrotas y diferencia media $+1.00$; esta sensibilidad no sustituye la tabla íntegra. No se aplica inferencia estadística a esta muestra pequeña y de condiciones desconocidas. Los cinco reportes externos se mantienen separados de las 19 partidas internas y no completan el plan de 104.

\textit{Guía: 5. Análisis Crítico Requerido; 7. Descripción de Funciones en el Artículo.}
Los marcadores muestran que Ispis perdió cuatro de cinco encuentros, pero sin las jugadas completas ni la configuración rival no puede identificarse la causa exacta. Como hipótesis de trabajo, la búsqueda a profundidad 4 puede omitir consecuencias más lejanas y la evaluación fija no distingue estabilidad real ni riesgo contextual de conceder esquinas; estas limitaciones podrían influir, pero las capturas no lo demuestran. Para mejorar, conviene registrar y reproducir cada jugada, localizar dónde se deterioran movilidad y acceso a esquinas, contrastar pesos y penalizaciones de casillas adyacentes a esquinas vacías con partidas de prueba separadas, y ensayar mayor profundidad dentro del mismo presupuesto de tiempo. Cada cambio debe evaluarse por separado con aperturas variadas y ambos colores; podría aumentar la competitividad, sin garantizar derrotar a esos rivales.
